\chapter{Atividades acadêmicas}

\section{Monitoria}
    \par Pode-se entender como monitor, um discente que juntamente com um docente orienta grupos de estudantes em atividades acadêmicas. Funcionando como um instrutor auxiliar do professor.  \par Há na UFBa, um projeto que visa através da monitoria “intensificar e assegurar a cooperação entre estudantes e professores nas atividades básicas da Universidade, relativas ao ensino, à pesquisa e à extensão, assim como subsidiar trabalhos acadêmicos, orientados por professores, através de ações multiplicadoras.” (PROGRAD UFBa).  \par São incluídas no projeto duas modalidades, monitoria com bolsa e monitoria voluntária, ambas com a garantia de certificado. Para inscrição no projeto deve-se preencher um formulário contido no site e o docente responsável deve apresentar a requisição junto ao seu departamento.

\section{Iniciação Científica, em Extensão e Pesquisa} \index{CR}
    \par A universidade não se limita apenas à área de ensino com aulas e provas, ela vai muito além disso com a parte de pesquisa e extensão, que é muito importante para complementar a formação dos estudantes e todos devem aproveitem muito essas atividades. Além disso, a participação em atividades de pesquisa e extensão é um dos requisitos para a conclusão do curso e permite o surgimento de outras oportunidades, tanto no ambiente interno, como no externo. A seguir, são apresentados os grupos de pesquisa e extensão do departamento de ciência da computação.
    
    \subsection{GRACO}
        \par Os Gestores da Rede Acadêmica de Computação (GRACO), é um grupo formado por estudantes do departamento de ciência da computação que tem como objetivo dar manutenção e aprimorar os serviços da rede do DCC. Participando do Graco é possível adquirir conhecimento em Sistemas Operacionais e administração de redes.
        
    \subsection{Grupos de Pesquisa}
    
        \subsubsection{Laboratório de Sistemas Distribuídos}
        
        \begin{wrapfigure}{L}{0.3\textwidth}
            \centering
            \includegraphics[width=0.29\textwidth]{lasid.png}
            %\caption{\label{logo:lasid}LaSid - Laboratório de Sistemas Distribuídos}
        \end{wrapfigure}
        O Laboratório de Sistemas Distribuídos (LaSiD) atua no desenvolvimento de técnicas e ferramentas para concepção dos sistemas distribuídos. O grupo foi criado no ano de 1995 e se tornou uma importante referência na América Latina, na área que atua. As atividades do laboratório se estendem também para a área de algoritmos distribuídos, redes de sensores, plataformas de software, entre outros.
        
        \subsubsection{Laboratório de Engenharia de Software}
        
        \begin{wrapfigure}{R}{0.3\textwidth}
            \centering
            \includegraphics[width=0.29\textwidth]{les.jpg}
            %\caption{\label{logo:les}LES - Laboratório de Engenharia de Software.}
        \end{wrapfigure}
        \par LES é a sigla para o Laboratório de engenharia de software, que tem como objetivo estudar a disciplina de engenharia de software e algumas subáreas, como, por exemplo, a área de design de software, sistemas sensíveis a contexto, métodos formais, reuso de software, entre outros. 
        
        \par O FORMAS é um grupo que atua na área de web semântica e ontologias de domínio e abrange alunos tanto da graduação como também da pós-graduação. O grupo iniciou suas atividades no ano de 2009 e vem realizando uma série de minicursos e palestras para divulgação da Web semântica. Os interessados devem procurar a professora Daniela Claro.
     
        \par Existem ainda outros grupos de pesquisa e extensão no departamento de ciência da computação como o grupo de pesquisa em redes de alto desempenho (GRADE), que atualmente desenvolve os projetos OpenWIMesh, que busca implantar uma infraestrutura de comunicação sem fio e o OpenWiMesh-MOB, projeto de estratégia de mobilidade em Redes Mesh sem fio definidas por software. Há também o Gaudi, grupo de algoritmos e computação distribuída, o WISER, que atua na área de Iternet, Web semântica, sistemas de recomendação da Web, computação em nuvens, entre outras áreas, o CEManTIKA, que estuda conceitos, técnicas, modelos e ferramentas que apoiem a construção de sensíveis ao contexto.
    
    \subsection{Onda Digital}
        \par É um programa de pesquisa, extensão e ensino com o objetivo de promover ações educativas, sociais e culturais, através da realização de cursos na área de informática e manutenção de equipamentos como uma forma de inclusão digital de um público que esteja em uma situação de exclusão social.
        
        \begin{wrapfigure}{R}{0.3\textwidth}
            \centering
            \includegraphics[width=0.29\textwidth]{gr.png}
            %\caption{\label{logo:lasid}LaSid - Laboratório de Sistemas Distribuídos}
        \end{wrapfigure}
    \subsection{GRUPRO}
        \par O Grupo de Programação (GRUPRO) busca aprimorar as habilidades em programação dos estudantes e para isso são realizados treinamentos sempre aos sábados a partir das 13:00 no laboratório 143 do instituto de Matemática, além de participarem de várias competições na área de programação como a maratona de programação e a olimpíada brasileira de informática. Interessados devem procurar os professores Maurício Pamplona ou Rubisley Lemes.
        
    
\section{Olimpíadas Universitárias}
\begin{wrapfigure}{R}{0.25\textwidth}
            \centering
            \includegraphics[width=0.24\textwidth]{ob}
\end{wrapfigure}
    \subsection{OBI}
        \par A Olímpiada Brasileira de Informática (OBI) é promovida pela Sociedade Brasileira de Computação e tem várias modalidades que vão desde o ensino fundamental até o nível universitário. Na modalidade Universitária podem competir somente os alunos que estiverem no primeiro ano da graduação.
        
    \subsection{Maratona de Programação}
    
        \begin{wrapfigure}{L}{0.25\textwidth}
            \centering
            \includegraphics[width=0.24\textwidth]{prog.jpg}
            %\caption{\label{logo:lasid}LaSid - Laboratório de Sistemas Distribuídos}
        \end{wrapfigure}
        \par Outra atividade acadêmica em que os alunos UFBA podem participar, é a maratona de programação, que é um evento do Associação para Máquinas da Computação (ACM) em associação com Sociedade Brasileira de Computação, que completará 20 anos de existência este ano. Ela é voltada para alunos dos cursos das áreas ligadas à computação, podendo está cursando a graduação ou início da pós-graduação. ``Os times são compostos por três alunos, que tentarão resolver durante 5 horas o maior número possível dos 8 ou mais problemas que são entregues no início da competição.'' (Maratona IME, USP).
        
    \subsection{ICPC}
        \par O International Collegiate Programing Contest (ICPC) é a etapa mundial da Maratona de Programação, e conseguem participar só as equipes que tiverem as melhores classificações na etapa nacional da Maratona de programação. É organizado pela Associação para Máquinas da Computação (ACM) e as provas são sempre em inglês. A próxima competição do ICPC vai acontecer na Dakota do Sul nos EUA.
       
    \subsection{IEEExtreme 24-Hour}
        \par É uma competição global de programação onde estudantes se confrontam em resoluções de problemas durante o período de 24 horas.
        \begin{figure}[!htb]
            \centering
            \includegraphics[scale = 0.3]{ie.jpg}
        \end{figure}

\begin{wrapfigure}{R}{0.3\textwidth}
    \centering
    \includegraphics[scale = 0.4]{pe.jpg}
\end{wrapfigure}
\section{Programa de Educação Tutorial}
    \par É um programa da Secretaria de Educação Superior, juntamente com o MEC que visa apoiar atividades acadêmicas que integrem a combinação de ensino, pesquisa e extensão. O objetivo central é formar grupos de alunos que trabalharão em conjunto com um professor-tutor, onde deverão realizar atividades extracurriculares de ensino, pesquisa e extensão, complementando sua formação acadêmica.
    \par De acordo com o ProGrad, os pets em atual funcionamento são:
    \begin{itemize}
        \item[$\star$] PET Comunicação;
        \item[$\star$] PET Enfermagem;
        \item[$\star$] PET Filosofia;
        \item[$\star$] PET Letras;
        \item[$\star$] PET Medicina;
        \item[$\star$] PET Odontologia;
        \item[$\star$] PET Comunidades Indígenas;
        \item[$\star$] PET Comunidades Populares;
        \item[$\star$] PET Eng. Sanitária e Ambiental;
        \item[$\star$] PET Nutrição;
        \item[$\star$] PET Pedagogia;
        \item[$\star$] PET Bacharelado Interdisciplinar;
        \item[$\star$] PET Farmácia - IMS.
    \end{itemize}

\section{Empresas Juniores}

    \begin{wrapfigure}{L}{0.3\textwidth}
        \centering
        \includegraphics[width=0.29\textwidth]{inf.png}
        %\caption{\label{logo:lasid}LaSid - Laboratório de Sistemas Distribuídos}
    \end{wrapfigure}
    \par As empresas juniores são associações sem fins lucrativos cujo principal objetivo é pôr em prática aquilo que foi aprendido na sala de aula, como também aprender conteúdos e ferramentas que não vistas na sala de aula e que são muito utilizadas no mercado de trabalho, além de possibilitar o ganho de experiência empresarial. A InfoJr é a empresa Júnior do departamento de Ciência da Computação, que foi criada no ano de 1998 e atua na área de desenvolvimento web, utilizando ferramentas bastante utilizados no mercado de trabalho como o HTML, CSS, PHP, JavaScript, MySQL, GIT e GITHUB. 
    
    \par A empresa realiza um processo de seleção a cada semestre para alunos dos cursos de Bacharelado em Ciência da Computação, Licenciatura em Ciência da Computação e Sistemas da Informação que estejam interessados em ingressar na empresa. Os alunos aprovados no processo de seleção passam por um processo Trainee que tem duração de 4 meses, onde são realizadas capacitações e os alunos devem desenvolver um projeto determinado pelos membros efetivos.
  

\section{Congressos}
    \par Um congresso é uma reunião entre pessoas que compartilham um mesmo interesse, visando tratar de assuntos relacionados a trabalhos científicos, apresentar propostas, trocar conhecimento, entre outros. Muitos congressos são realizados com participação de alunos universitários, a seguir são listados alguns exemplos de congressos que os estudantes podem participar:
        
\subsection{Fórum Mundial}
        \par O fórum mundial é um congresso realizado anualmente e é separado em vários tipos, como o fórum mundial da agua, fórum mundial das migrações, fórum mundial da educação profissional e tecnológica, etc. \par Estes fóruns são realizados em uma cidade brasileira predeterminada e é possível a participação de qualquer pessoa. Para participar como palestrante é necessário submeter um projeto, relacionado com o tema escolhido pelo fórum no ano em questão, à seleção interna do evento. Já para participar como visitante é apenas necessário se inscrever pelo site do evento.
\subsection{Congresso UFBA}
\begin{wrapfigure}{R}{0.3\textwidth}
            \centering
            \includegraphics[scale = 0.4]{cong.png}
        \end{wrapfigure}
        O congresso UFBA 70 anos foi um evento realizado no ano de 2016, do dia 14/07 a 17/07, que visava uma reflexão sobre a universidade, sobre as consequências dela na vida de seus alunos e ex-alunos, suas perspectivas para o futuro e seu papel social. \par Neste congresso houveram participações de ex-alunos, alunos, pessoas interessadas nos assuntos tratados e profissionais de áreas diversas, com interesse nas palestras, oficinas, discussões e minicursos realizados durante os dias do evento.
        
        
    \subsection{SEMCOMP}
        \par A SEMCOMP ou Semana de Computação é um evento realizado anualmente na UFBA, na escola politécnica, e conta com mais de 30 palestras, com participação de representantes de grandes empresas e desenvolvedores como palestrantes, e 4 cursos, com temas relacionados com a área de computação. \par Este evento é importante para trazer uma visão mais ampla da atuação de profissionais da área de TI, bem como criar um possível contato entre os profissionais da área e os estudantes.
